\documentclass[11pt,a4paper,oneside]{book}
\usepackage[margin=2.6cm,headheight=14pt]{geometry}
\usepackage{cmbright}
\usepackage[T1]{fontenc}
\renewcommand{\ttdefault}{pcr}
\usepackage{fancyhdr}
\usepackage[nohints]{minitoc}
\usepackage{amsmath,amssymb,bm}
\usepackage{graphicx}
\usepackage{booktabs}
\usepackage{array}
\usepackage{algorithm}
\usepackage{algpseudocode}
\usepackage[numbers,sort&compress]{natbib}
\usepackage[colorlinks=true,linkcolor=blue!50!black,citecolor=green!40!black,urlcolor=blue!60!black]{hyperref}
\usepackage{xcolor}

\newcommand{\vth}{\bm\theta}
\newcommand{\vout}{\bm v_{\mathrm{out}}}
\newcommand{\onec}{\mathrm{1c}}
\newcommand{\pbezero}{PBE0r}
\newcommand{\scanzero}{SCAN0r}
\newcommand{\pbeslzero}{PBE-SL0r}
\newcommand{\key}[1]{\textbf{\texttt{#1}}}
\setcounter{minitocdepth}{2}
\setlength{\mtcindent}{12pt}
\renewcommand{\mtctitle}{}
\mtcsetrules{minitoc}{off}
\mtcsetfont{minitoc}{section}{\normalfont\small}
\mtcsetfont{minitoc}{subsection}{\normalfont\small}
\mtcsettitlefont{minitoc}{\normalfont}
\pagestyle{fancy}
\fancyhf{}
\renewcommand{\chaptermark}[1]{\markboth{\thechapter\ \MakeUppercase{#1}}{}}
\fancyhead[R]{\small\leftmark\quad\thepage}
\renewcommand{\headrulewidth}{0.4pt}
\fancypagestyle{plain}{\fancyhf{}\fancyfoot[C]{\small\textbf{\thepage}}\renewcommand{\headrulewidth}{0pt}}
\newcommand{\bra}[1]{\langle #1|}
\newcommand{\ket}[1]{|#1\rangle}
\newcommand{\braket}[2]{\langle #1|#2\rangle}
\newcommand{\pt}{\tilde p}

\title{\Huge\bfseries VESPAW\\[0.6em]\LARGE Theory and methods\\[2em]
\normalsize Variational electronic structure with projector augmented waves}
\author{Mohsen Sotoudeh}
\date{Version 1.0.2}

\begin{document}
\dominitoc
\frontmatter
\maketitle
\tableofcontents

\mainmatter
% =====================================================================
\chapter{Introduction}
\minitoc
% =====================================================================

\section{Why another look at an old problem}

A first-principles calculation of a material ends with a number, a
total energy, a band gap, a magnetic moment. Whether this number means
anything depends on a long chain of steps: the density functional, the
basis set, the treatment of the atomic cores, the sampling of the
Brillouin zone, and finally the iteration that makes the electronic
structure self-consistent. Most of these steps are controlled by a few
numbers that can be converged. The last one is different. A
self-consistency cycle that has stalled, or that has settled in a
metastable state, may report convergence just like one that has found
the ground state, and nothing in the output tells the difference.

For an experienced user who looks at every calculation, this is a
nuisance. For the large sets of calculations that are now used to
search for new materials for batteries, catalysts and other energy
technologies, it is a real problem, because nobody looks at every
calculation any more. The materials that matter most, transition-metal
oxides with magnetic order, mixed valence and polarons, are exactly
the ones where self-consistency is most delicate.

VESPAW starts from this observation. Its central idea is to treat the
effective Kohn--Sham Hamiltonian, and not the density, as the variable
of the self-consistency problem, and to evaluate the energy for the
actual state that each trial Hamiltonian produces. Every energy is then
an upper bound of the ground-state energy, and the iteration accepts a
step only if this bound does not increase. In this book I describe the
physics and the numerical methods on which VESPAW rests, with the
emphasis on how the Kohn--Sham equations are solved.

\section{VESPAW and CP-PAW}

VESPAW is a plane-wave code based on the projector augmented-wave (PAW)
method~\cite{bloechl94_prb50_17953}. It is a fork of CP-PAW, the
Car--Parrinello PAW code written by Peter E.~Bl\"ochl at Clausthal
University of Technology. The PAW formalism, the construction of the
setups, the Car--Parrinello dynamics, the local orbitals and the local
hybrid functional \pbezero, and the input format all come from CP-PAW.
VESPAW adds
\begin{itemize}
\item a \emph{variational} self-consistency optimizer, whose energies
  are upper bounds of the ground-state energy in every iteration
  (chapter~\ref{chap:variational});
\item iterative eigensolvers (block Davidson, RMM-DIIS, conjugate
  gradients) next to the damped dynamics of CP-PAW
  (chapter~\ref{chap:eigensolvers});
\item functionals checked against libxc and VASP, the SCAN family
  (SCAN-L, and SCAN and r$^2$SCAN with the orbital kinetic-energy
  density), and the
  local hybrids \pbeslzero\ and \scanzero, which build on \pbezero\
  (chapters~\ref{chap:dft} and~\ref{chap:hybrids});
\item quality checks for PAW setups, an audited setup library
  (chapter~\ref{chap:paw}) and a regression suite
  (chapter~\ref{chap:code}).
\end{itemize}

\paragraph{License.} CP-PAW is distributed under the GNU General Public
License, version~3 (GPL-3.0). VESPAW is derived from CP-PAW and is
distributed under the same license. The older notice ``(C) Copyright
Clausthal University of Technology, licensed materials'' in the headers
of the inherited source files predates the GPL release.

\section{How to read this book}

Chapters~\ref{chap:dft} to~\ref{chap:paw} summarize the physics that
VESPAW implements: density-functional theory, plane waves and the PAW
method. Chapter~\ref{chap:hybrids} treats the local hybrid functionals.
Chapters~\ref{chap:variational} and~\ref{chap:eigensolvers} describe how
the Kohn--Sham equations are solved, which is where VESPAW differs most
from other codes. Chapter~\ref{chap:forces} covers forces, stress and
molecular dynamics, and chapter~\ref{chap:code} the structure of the
program. The hands-on tutorial (\texttt{doc/tutorial.pdf}) puts all of
this into practice.

Input keywords are written in \key{bold typewriter}; their complete
description is in the input manual. Hartree atomic units are used
throughout unless stated otherwise, while cutoff energies are given in
Rydberg, as in the input. The program reports its results in eV, \AA{}
and GPa.

\paragraph{Pitfalls.} Many sections end with a short list of pitfalls.
They record errors that were actually found in the code and fixed. I
have kept them, because they show which quantities are delicate and how
they are tested.

% =====================================================================
\chapter{Density-functional theory}\label{chap:dft}
\minitoc
\section{Introduction}

The electrons of a solid interact with each other, and their wave
function depends on the coordinates of all of them at once. Already
for a few dozen electrons such a wave function cannot even be stored,
let alone calculated. Density-functional theory removes this obstacle
in two steps. Hohenberg and Kohn showed that the ground-state energy is
a functional of the electron density alone, a function of only three
coordinates. Kohn and Sham then expressed this functional through a
fictitious system of non-interacting electrons that has the same
density. All the many-body physics is collected in one term, the
exchange-correlation energy, and the art of density-functional theory
lies in its approximation. This chapter summarizes the theory and the
functionals implemented in VESPAW, with particular attention to how
their parameters were checked against independent implementations.

% =====================================================================

\section{Kohn--Sham equations}

By the Hohenberg--Kohn theorem~\cite{HohenbergKohn1964} the ground-state
energy of $N$ interacting electrons in an external potential
$v_{\mathrm{ext}}$ is a functional of the density $n(\bm r)$. Kohn and
Sham~\cite{KohnSham1965} write it in terms of non-interacting orbitals
$\psi_n$ with occupations $f_n$,
\begin{equation}
  E[n] = \sum_n f_n \bra{\psi_n}-\tfrac12\nabla^2\ket{\psi_n}
  + \int d^3r\, v_{\mathrm{ext}} n
  + \frac12\iint d^3r\,d^3r'\,\frac{n(\bm r)n(\bm r')}{|\bm r-\bm r'|}
  + E_{\mathrm{xc}}[n],
  \qquad n(\bm r)=\sum_n f_n|\psi_n(\bm r)|^2 .
\end{equation}
Minimizing with orthonormal orbitals gives the Kohn--Sham equations
\begin{equation}
  \Bigl[-\tfrac12\nabla^2 + v_{\mathrm{eff}}(\bm r)\Bigr]\psi_n = \epsilon_n\psi_n,
  \qquad
  v_{\mathrm{eff}} = v_{\mathrm{ext}} + v_{\mathrm H}[n] + \frac{\delta E_{\mathrm{xc}}}{\delta n}.
\end{equation}
Because $v_{\mathrm{eff}}$ depends on the density, the equations must be
solved self-consistently (Chapter~\ref{chap:variational}). A
self-consistent solution is a stationary point of $E$ only if the
potential is the functional derivative of the energy. This condition is
easy to state and easy to violate in an implementation; VESPAW tests it
routinely (Sec.~\ref{sec:consistency}).

\section{Exchange-correlation functionals}

\paragraph{LDA and GGA.} The local-density approximation uses the
exchange-correlation energy per electron of the homogeneous electron gas,
$E_{\mathrm{xc}}=\int n\,\epsilon_{\mathrm{xc}}(n)$ (PZ~\cite{perdew81_prb23_5048},
PW92~\cite{perdew92_prb45_13244}). Generalized gradient approximations add a
dependence on $|\nabla n|$; the default in VESPAW is
PBE~\cite{perdew96_prl77_3865}. All functionals are selected by
\key{!CONTROL!DFT TYPE}, as a number or a name, e.g.\ \key{TYPE='PBE'};
the VASP \key{GGA} tags are accepted as names. Table~\ref{tab:functionals}
lists them.

\begin{table}[tb]
\centering\small
\caption{Exchange-correlation functionals of VESPAW. ``libxc'' marks
functionals whose energy density agrees with libxc~7.1.2 (and the VASP
implementation) to $\sim10^{-15}$ at every point of a grid in density,
gradient, spin polarization and $\alpha$ (October 2026). The local
hybrids are described in Chapter~\ref{chap:hybrids}.}
\label{tab:functionals}
\begin{tabular}{@{}r>{\raggedright\arraybackslash}p{2.3cm}>{\raggedright\arraybackslash}p{6.6cm}lc@{}}
\toprule
\key{TYPE} & name & functional & VASP & libxc \\
\midrule
\multicolumn{5}{@{}l}{\emph{LDA}}\\
1  & \key{PZ}   & Perdew--Zunger~\cite{perdew81_prb23_5048} of Ceperley--Alder~\cite{ceperley80_prl45_566} & & \checkmark\\
2  & \key{PW92} & Perdew--Wang 92~\cite{perdew92_prb45_13244} & & \checkmark\\
3  &            & Slater exchange only & & \checkmark\\
4  &            & X$_\alpha$, $\alpha=0.7$ & & \\
\multicolumn{5}{@{}l}{\emph{GGA}}\\
10 & \key{PBE}, \key{PE}    & PBE~\cite{perdew96_prl77_3865} & \key{GGA=PE} & \checkmark\\
11 & \key{RPBE}, \key{RP}   & RPBE~\cite{hammer99_prb59_7413} & \key{GGA=RP} & \checkmark\\
12 & \key{PW91}, \key{91}   & PW91 exchange and correlation~\cite{perdew92_prb46_6671} & \key{GGA=91} & \checkmark\\
13 & \key{REVPBE}, \key{RE} & revPBE, $\kappa=1.245$~\cite{zhang98_prl80_890} & \key{GGA=RE} & \checkmark\\
14 & \key{PBESOL}, \key{PS} & PBEsol~\cite{perdew08_prl100_136406} & \key{GGA=PS} & \checkmark\\
5001 & \key{BLYP} & B88~\cite{becke88_pra38_3098} + LYP~\cite{lee88_prb37_785} & & \checkmark\\
81 &  & BP86: PZ + B88 + P86~\cite{perdew86_prb33_8822} & & \checkmark\\
8  &  & B88 + P86 on PW92 local correlation (non-standard) & & \\
7, 71, 9 & & B88 exchange with PW92, PZ or PBE correlation & & \checkmark\\
6  &  & B88 exchange only & & \checkmark\\
\multicolumn{5}{@{}l}{\emph{meta-GGA}}\\
21 & \key{SCAN} & SCAN, orbital $\tau$~\cite{sun15_prl115_036402} (with \key{TSCANDELOC=T}) & \key{METAGGA=SCAN} & \checkmark\\
22 & \key{R2SCAN} & r$^2$SCAN, orbital $\tau$~\cite{furness20_jpcl11_8208} (with \key{TSCANDELOC=T}) & \key{METAGGA=R2SCAN} & \checkmark\\
20 & \key{SCANL} & SCAN-L, deorbitalized~\cite{mejiarodriguez17_pra96_052512,perdew07_prb75_155109} & \key{METAGGA=SCANL} & \checkmark\\
\multicolumn{5}{@{}l}{\emph{local hybrids} (Chapter~\ref{chap:hybrids})}\\
   & \pbezero   & PBE + onsite exact exchange & & \\
   & \pbeslzero & PBE; onsite SCAN-L/exact exchange & & \\
   & \scanzero  & SCAN + onsite exact exchange & & \\
   & r$^2$SCAN0r  & r$^2$SCAN + onsite exact exchange & & \\
\multicolumn{5}{@{}l}{\emph{for testing and reproducing old results}}\\
$-10$ & & PBE before the 2024/2026 corrections (see text) & & \\
$-12$ & & PBE exchange + PW91 correlation (old type 12) & & \\
5002 & & PBE with exchange scaled by 0.75 & & \\
10001--7 & & single components (PZ, PW92, P86, PBE, PW91 correlation; PBE, RPBE exchange) & & \\
\bottomrule
\end{tabular}
\end{table}

\paragraph{Parameters.} The constants follow the original publications
and agree with libxc and VASP: $\beta=0.06672455060314922$ in PBE
correlation and in $\mu=\beta\pi^2/3$ of PBE exchange (before October 2026
VESPAW used $\beta=\nu C_c(0)=0.0667263$, kept by \key{TYPE=-10} together
with the pre-2024 reduced gradient, $t^2$ too small by a factor 0.62);
PW91 keeps $\beta=\nu C_c(0)$, as in VASP; PBEsol uses $\mu=10/81$ and
$\beta=0.046$; SCAN-L uses the optimized Perdew--Constantin parameters
$a=1.784720$, $b=0.258304$ of Ref.~\cite{mejiarodriguez17_pra96_052512};
the spin interpolation of SCAN correlation uses
$G_c(\zeta)=[1-2.363(d_x(\zeta)-1)](1-\zeta^{12})$, with 2.363 as in VASP
and libxc (the SCAN supplement quotes 2.3631). Unit conversions use
CODATA~2018 (1~Hartree $=27.211386$~eV, 1~bohr $=0.529177211$~\AA).
Plane-wave hybrids with full exact exchange (PBE0, HSE06, SCAN0, B3LYP)
are not yet available.

\paragraph{Meta-GGAs: SCAN, r$^2$SCAN, SCAN-L and TSCAN.} The SCAN
functional~\cite{sun15_prl115_036402} depends in addition on the kinetic
energy density $\tau=\frac12\sum_n f_n|\nabla\psi_n|^2$ through the
iso-orbital indicator
\begin{equation}
  \alpha = \frac{\tau-\tau_W}{\tau_{\mathrm{unif}}},\qquad
  \tau_W=\frac{|\nabla n|^2}{8n},\qquad
  \tau_{\mathrm{unif}}=\tfrac{3}{10}(3\pi^2)^{2/3}n^{5/3}.
\end{equation}
VESPAW provides two routes:
\begin{itemize}
\item \textbf{SCAN-L} (\key{TYPE='SCANL'}), the deorbitalized SCAN of
  Mej\'{\i}a-Rodr\'{\i}guez and Trickey~\cite{mejiarodriguez17_pra96_052512,mejiarodriguez18_prb98_115161},
  in which $\alpha$ is approximated from $n$, $|\nabla n|$ and $\nabla^2 n$.
  The functional depends on the Laplacian of the density, so its potential
  contains second-derivative terms (see below).
\item \textbf{TSCAN} (\key{TYPE='SCAN'} with \key{TSCANDELOC=T}), the
  orbital-dependent SCAN with the true $\tau$. The Hamiltonian then
  contains the operator $-\frac12\nabla\cdot v_\tau\nabla$ with
  $v_\tau=\partial e_{\mathrm{xc}}/\partial\tau$. In the PAW spheres the
  one-center $\tau$ is built from the partial waves, including its
  non-spherical parts.
\item \textbf{r$^2$SCAN} (\key{TYPE='R2SCAN'} with \key{TSCANDELOC=T}),
  the regularized and restored SCAN of Furness \emph{et
  al.}~\cite{furness20_jpcl11_8208}, with the same orbital $\tau$ and the
  same PAW treatment as TSCAN. It replaces $\alpha$ by the regularized
  $\bar\alpha=(\tau-\tau_W)/(\tau_{\mathrm{unif}}+\eta\tau_W)$ with
  $\eta=10^{-3}$, the switching function of SCAN by a polynomial for
  $0\le\bar\alpha\le2.5$, and adds the gradient-expansion corrections of
  exchange and correlation. The implementation follows the paper and its
  erratum and agrees with libxc to $10^{-14}$. Band gaps change smoothly
  with the cutoff, unlike those of SCAN, which scatter by a few 0.01~eV.
\end{itemize}

\paragraph{PAW treatment of $\tau$.} The kinetic-energy density is
divided like the density: the plane-wave part $\tilde\tau$ of the
auxiliary wave functions plus the pseudized core kinetic-energy density
$\tilde\tau_c$, and in each sphere the one-center terms
$\tau^1-\tilde\tau^1=\sum_{ij}D_{ij}\bigl(\tfrac12\nabla\phi_i\cdot\nabla\phi_j
-\tfrac12\nabla\tilde\phi_i\cdot\nabla\tilde\phi_j\bigr)$ plus the core
parts. Two checks are useful: $\int\tilde\tau\,d^3r$ equals the kinetic
energy of the auxiliary wave functions, and
$\int(\tau^1-\tilde\tau^1)\,d^3r=\sum_{ij}D_{ij}\Delta T_{ij}$. The setups
of the meta-GGAs are constructed with PBE, as for VASP.

\paragraph{Functional derivatives.} For a semilocal energy
$E=\int e(n,|\nabla n|^2,\nabla^2 n)\,d^3r$ the potential is the full
Euler--Lagrange derivative
\begin{equation}
  v = \frac{\partial e}{\partial n}
  - \nabla\cdot\Bigl(2\frac{\partial e}{\partial|\nabla n|^2}\nabla n\Bigr)
  + \nabla^2\frac{\partial e}{\partial \nabla^2 n}.
  \label{eq:euler}
\end{equation}
For GGAs the last term is absent; for SCAN-L it is essential. Where the
energy carries a spatial weight (the double counting of the local
hybrids, Sec.~\ref{sec:pbesl0r}), the weight stands \emph{inside} the
derivatives. In the spherical one-center expansion the divergence
becomes $r^{-2}(r^2 A)'$; in the non-spherical case the weak form is
used, projecting $\partial e/\partial\nabla n$ onto the radial and
tangential gradients of each spherical harmonic.

\paragraph{Pitfall.} Before September 2026, the one-center SCAN-L and
TSCAN potentials kept only the first term of Eq.~\eqref{eq:euler}.
Self-consistent solutions were then not stationary points of the energy:
the variational optimizer found lower energies without reaching
self-consistency, and for TSCAN it did not converge at all. The
consistency test of Sec.~\ref{sec:consistency} exposed both.

\paragraph{Pitfall (October 2026).} Three errors made all meta-GGA band
gaps too small by about 0.3~eV, e.g. 4.93 instead of 5.24~eV for SCAN MgO:
the setups were built from an atom without exchange and correlation, the
one-center $\tau$ lacked the factor $\frac12$, and the pseudized core
$\tau$ was missing in the plane-wave part. The energy and the
Hamiltonian were consistent with each other, so the consistency tests
passed. They were found with the two checks above and by comparison
with VASP: SCAN and r$^2$SCAN gaps now agree within 0.06~eV (MgO 5.30
and 5.34~eV, VASP 5.24 and 5.30~eV; Si 0.949 and 0.919~eV, VASP 0.969
and 0.917~eV). For SCAN0r the double counting also subtracted the
potential of the core $\tau$, which made the Hamiltonian inconsistent
with the energy and stalled the variational optimizer.

\section{Spin}

Collinear spin polarization (\key{!STRUCTURE!OCCUPATIONS NSPIN=2}) uses two
sets of orbitals and the spin densities $n_\uparrow,n_\downarrow$; the
code works internally with total and spin density $(n,m)$.
Non-collinear magnetism uses two-component spinors. For exchange, the
spin-scaling relation
\begin{equation}
  E_x[n_\uparrow,n_\downarrow]=\tfrac12\bigl(E_x[2n_\uparrow]+E_x[2n_\downarrow]\bigr)
  \label{eq:spinscaling}
\end{equation}
is exact and allows any spin-unpolarized exchange functional to be
applied to a spin-polarized density. VESPAW uses it for the exchange of
the SCAN family, on the grid and in the one-center terms. Correlation
does not obey Eq.~\eqref{eq:spinscaling}; SCAN correlation is evaluated
for the spin densities with the spin-interpolation of
Ref.~\cite{sun15_prl115_036402}: $\phi(\zeta)$ in the gradient term,
$G_c(\zeta)$ in the single-orbital limit, and $\alpha$ divided by
$d_s(\zeta)=[(1+\zeta)^{5/3}+(1-\zeta)^{5/3}]/2$ with the total $\tau$.
For SCAN-L, the spin-resolved kinetic energy densities are deorbitalized
separately, $\tau[n_\uparrow,n_\downarrow]=\frac12(\tau[2n_\uparrow]+\tau[2n_\downarrow])$,
as in libxc.

\section{Finite temperature}

For metals and for systems with near-degenerate states, VESPAW minimizes
the Mermin free energy~\cite{Mermin1965}
\begin{equation}
  A = E[\{\psi_n\},\{f_n\}] - T S[\{f_n\}],\qquad
  S=-k_B\sum_n\bigl[f_n\ln f_n+(1-f_n)\ln(1-f_n)\bigr],
  \label{eq:mermin}
\end{equation}
with Fermi occupations $f_n=[1+e^{(\epsilon_n-\mu)/k_BT}]^{-1}$ at fixed
electron number. The electronic temperature is set in \key{!CONTROL!MERMIN
T[K]}. It acts as a smearing parameter; the free energy $A$, not $E$, is
the variational quantity.


\section*{Further reading}
The original work on the subject of this chapter is found in
Refs.~\cite{HohenbergKohn1964,KohnSham1965,Mermin1965,perdew96_prl77_3865,sun15_prl115_036402}.

% =====================================================================
\chapter{Plane waves and Brillouin-zone sampling}\label{chap:planewaves}
\minitoc
\section{Introduction}

Plane waves are the simplest basis set one can imagine. They are
orthonormal, they do not depend on the positions of the atoms, and
their completeness is controlled by a single number, the cutoff. The
fast Fourier transform switches between real and reciprocal space at
low cost, so that each term of the Hamiltonian can be applied where it
is simplest. The price is that sharp features, such as the oscillations
of the wave functions near a nucleus, would require an enormous number
of plane waves. The PAW method of the next chapter removes this
problem. Here we collect what the plane-wave description of a periodic
solid involves: Bloch states, cutoffs, the grids for the density, and
the sampling of the Brillouin zone.

% =====================================================================

\section{Bloch states and plane waves}

In a crystal with lattice vectors $\bm T$ the orbitals are Bloch states,
$\psi_{n\bm k}(\bm r+\bm T)=e^{i\bm k\cdot\bm T}\psi_{n\bm k}(\bm r)$,
which are expanded in plane waves,
\begin{equation}
  \tilde\psi_{n\bm k}(\bm r)=\sum_{\bm G} c_{n\bm k}(\bm G)\,
  e^{i(\bm k+\bm G)\cdot\bm r},
  \qquad \tfrac12|\bm k+\bm G|^2\le E_{\mathrm{cut}} .
\end{equation}
In PAW only the smooth pseudo wave functions $\tilde\psi$ are expanded in
plane waves (Chapter~\ref{chap:paw}).

\section{Cutoffs and FFT grids}

The wave-function cutoff is \key{!CONTROL!FOURIER EPWPSI} in Rydberg. The
density contains products of wave functions and would need four times the
cutoff; VESPAW uses \key{EPWRHO = CDUAL$\times$EPWPSI}. \key{CDUAL=4} is
exact, \key{CDUAL=2} is usually sufficient and cheaper. Density and local
potential live on a real-space grid; the Hamiltonian is applied to a
wave function by transforming it to real space with a fast Fourier
transform, multiplying by the potential and transforming back. In VESPAW
these transforms are batched over blocks of 32 bands.

Energies converge from above with the cutoff, but energy
\emph{differences}, forces and stress converge faster. The required cutoff
is set by the hardest PAW setup in the system; the setup generator
prints an estimate (Sec.~\ref{sec:kintable}).

\section{$k$-point sampling}

Brillouin-zone integrals are replaced by sums over a Monkhorst--Pack
mesh~\cite{MonkhorstPack1976}, \key{!STRUCTURE!KPOINTS DIV=$n_1$ $n_2$
$n_3$}, with an optional shift, or by a target density \key{R}.
Symmetry and time reversal ($\psi_{-\bm k}=\psi_{\bm k}^*$) reduce the
number of irreducible points. Metals need dense meshes together with a
finite electronic temperature; insulators converge quickly.

\section{The $\Gamma$ point}

At $\bm k=0$ the wave functions can be chosen real. CP-PAW exploits this
by packing two real bands into one complex array (``super wave
functions''), which halves the number of FFTs. An alternative half-sphere
storage (\key{GAMMASUPER}) is experimental and not yet correct.


\section*{Further reading}
The original work on the subject of this chapter is found in
Refs.~\cite{Payne1992,MonkhorstPack1976}.

% =====================================================================
\chapter{The projector augmented-wave method}\label{chap:paw}
\minitoc
\section{Introduction}

A wave function in a solid has two faces. Between the atoms, where the
chemical bonds form, it is smooth. Close to a nucleus it oscillates
rapidly, because it must be orthogonal to the core states. The classic
electronic-structure methods have chosen one face. Pseudopotentials
replace the atomic region by a smooth effective potential and lose the
true wave function there. Augmented-wave methods keep it, but at the
price of a basis set that depends on energy and on the atomic
positions. The PAW method of Bl\"ochl unites both views through a linear
transformation between smooth auxiliary wave functions, which are
expanded in plane waves, and the true wave functions. This chapter
describes the transformation, the construction of the PAW setups, and
the checks that make a setup trustworthy.

% =====================================================================

\section{The PAW transformation}

The PAW method~\cite{bloechl94_prb50_17953} maps the true all-electron
(AE) valence orbitals $\psi_n$, which oscillate strongly near the nuclei,
onto smooth pseudo (PS) orbitals $\tilde\psi_n$ by a linear transformation
\begin{equation}
  \ket{\psi_n}=\ket{\tilde\psi_n}
  +\sum_{R,i}\bigl(\ket{\phi_{R,i}}-\ket{\tilde\phi_{R,i}}\bigr)
  \braket{\pt_{R,i}}{\tilde\psi_n}.
\end{equation}
Here $\phi_{R,i}$ are AE partial waves (solutions of the atomic
Schr\"odinger equation at reference energies), $\tilde\phi_{R,i}$ the
corresponding PS partial waves, which agree with $\phi_{R,i}$ outside an
augmentation radius $r_c$, and $\pt_{R,i}$ projector functions, localized
within $r_c$ and dual to the PS partial waves,
$\braket{\pt_{R,i}}{\tilde\phi_{R,j}}=\delta_{ij}$. The index $i$ collects
the angular momentum $(\ell,m)$ and a counter $\nu$ over partial waves of
the same $\ell$. The number of partial waves per $\ell$ is set by
\key{!STRUCTURE!SPECIES NPRO}.

Expectation values split into a plane-wave part and one-center
corrections,
\begin{equation}
  \langle A\rangle = \sum_n f_n\bra{\tilde\psi_n}\tilde A\ket{\tilde\psi_n}
  + \sum_R\sum_{ij}D_{R,ij}\bigl(\bra{\phi_{R,j}}A\ket{\phi_{R,i}}
  -\bra{\tilde\phi_{R,j}}\tilde A\ket{\tilde\phi_{R,i}}\bigr),
  \qquad
  D_{R,ij}=\sum_n f_n\braket{\pt_{R,i}}{\tilde\psi_n}\braket{\tilde\psi_n}{\pt_{R,j}},
\end{equation}
and the total energy becomes $E=\tilde E+\sum_R(E^1_R-\tilde E^1_R)$: a
plane-wave term plus, for each atom, the difference of an AE and a PS
one-center energy evaluated on radial grids. The orbitals are orthonormal
with respect to the overlap operator
\begin{equation}
  \hat O = 1+\sum_{R,ij}\ket{\pt_{R,i}}\Delta O_{R,ij}\bra{\pt_{R,j}},
  \qquad
  \Delta O_{R,ij}=\braket{\phi_{R,i}}{\phi_{R,j}}-\braket{\tilde\phi_{R,i}}{\tilde\phi_{R,j}},
\end{equation}
so the Kohn--Sham equations take the generalized form
$\tilde H\ket{\tilde\psi_n}=\epsilon_n\hat O\ket{\tilde\psi_n}$ with
\begin{equation}
  \tilde H=-\tfrac12\nabla^2+\tilde v(\bm r)
  +\sum_{R,ij}\ket{\pt_{R,i}}\Delta H_{R,ij}\bra{\pt_{R,j}} .
\end{equation}

\section{Compensation charges, core and local potential}

The electrostatics of the one-center densities is handled by
compensation charges $\hat n_R$ (Gaussians of width set by
\key{RCSM/RCOV}) that reproduce the multipole moments of
$n^1_R-\tilde n^1_R$, so that the long-range electrostatics is carried by
the smooth density $\tilde n+\hat n$ on the plane-wave grid. The
one-center expansions are truncated at $\ell\le$ \key{LRHOX}. Core
electrons are frozen in their atomic state. For the exchange-correlation
energy, which is non-linear in the density, a pseudized core density
(\key{!CORE POW RC}) is added to the smooth density. The local
pseudopotential is constructed by pseudizing the AE potential inside a
radius (\key{!POT POW RC}).

\section{Setup construction}\label{sec:setups}

A PAW \emph{setup} contains, for one element, the partial waves,
projectors, core densities and local potential. VESPAW builds setups at
run time from a short parameter block (\key{!SPECIES!AUGMENT}), or reads
them from the library \texttt{vespaw\_pot}. Three construction types exist:
\begin{itemize}
\item \key{HBS}: the scheme of Bl\"ochl's later CP-PAW versions. The
  first partial wave of each $\ell$ is the lowest valence state; further
  partial waves are obtained from the nodeless (``$\dot\phi$-like'')
  construction at higher energies, bounded by a box of radius
  \key{RBOX/RCOV}.
\item \key{NDLSS}: an older nodeless scheme.
\item \key{KERKER}: Kerker-type pseudization.
\end{itemize}
The radii are given in units of the covalent radius of the element:
\key{RCL/RCOV} for the partial waves, \key{RCSM/RCOV} for the compensation
charge, \key{RC/RCOV} in \key{!POT} and \key{!CORE}.

\paragraph{Semicore states.} When shallow core states are treated as
valence (e.g.\ Cu $3s3p$, \key{ZV=19}), an angular momentum has two
valence states, $3s$ and $4s$. Each valence bound state then receives its
own partial wave at its own energy; only further partial waves are box
states. With a semicore state below it, the PS $4s$ orbital need not have
a node, because it is orthogonal to the PS $3s$ only through $\hat O$. The
setup test therefore identifies PS states by their energy, not by their
node count.

\section{Quality checks}\label{sec:setupchecks}

A setup is an approximation, and a bad one fails silently: calculations
converge to wrong answers. The generator therefore reports three checks.

\paragraph{Ghost test.} The atomic PAW Hamiltonian of each $\ell$ is
diagonalized in a basis of 120 spherical Bessel functions inside a large
sphere. Its lowest levels must match the AE valence levels. A level far
below the lowest AE reference is a \emph{ghost state}; the generator
prints \texttt{WARNING: GHOST STATE IN SETUP}. The overlap matrix in this
basis must be positive definite; if it is not, the setup is unusable.

\paragraph{Incompleteness warning.} With only one partial wave per
$\ell$, the local potential may bind a state far below the valence level
whose wave function lies largely outside the span of the projectors.
Such components receive no one-center correction. In a solid they can
sink into the potential well of the compensation charge, and the
calculation converges to a collapsed state far below the physical
minimum. This happened for NiO with a single-projector Ni setup: the
collapse deepened with the cutoff and disappeared with a second $s$
projector. The generator prints \texttt{WARNING: INCOMPLETE SETUP} when
the lowest state of the local potential lies more than 0.5~H below the
valence reference with a single projector.

\paragraph{Kinetic-energy cutoff table.}\label{sec:kintable} As in the
VASP setup generator, the kinetic energy of each PS valence orbital that
lies above a cutoff is computed from its Fourier transform, weighted by
the occupation and summed. The generator prints the error as a function of
the cutoff, from 20~eV to about 290~Ry, and suggests \texttt{ENMAX} and
\texttt{ENMIN} (errors of 6 and 13~meV), or states that they are not
reached. The PS orbitals are tapered to zero before the box wall to avoid
a spurious kink.

\section{The setup library}

\texttt{vespaw\_pot} contains one species block per element and variant.
Its audit (\texttt{vespaw\_pot/AUDIT.md}) lists each setup with its
status (OK, incomplete, ghost, unusable) and the recommended replacement.
The flagged setups were rebuilt as HBS setups with two projectors per
$\ell$. Validation in solids shows that the frozen-semicore $3d$ setups
(Fe, Cu) reproduce AE-PBE lattice constants only to about 1\%, and that
the error grows with the augmentation radius; a semicore Cu setup with
small radii reaches 0.5\% at a high cutoff. Accurate soft $3d$ setups are
an open problem (see the audit).


\section*{Further reading}
The original work on the subject of this chapter is found in
Refs.~\cite{bloechl94_prb50_17953,KresseJoubert1999}.

% =====================================================================
\chapter{Local hybrid functionals}\label{chap:hybrids}
\minitoc
\section*{Introduction}
% =====================================================================

Hybrid functionals admix a fraction $a_x$ of exact exchange and cure
much of the band-gap problem of semilocal functionals, but in a
plane-wave basis exact exchange is one to two orders of magnitude more
expensive than a GGA. Local hybrids of the \pbezero\
type~\cite{bloechl11_prb84_205101,sotoudeh17_prb95_235150} keep only the
onsite part of the exchange correction, evaluated in a basis of local
orbitals, at nearly semilocal cost.

\section{Local orbitals}

Each atom carries natural tight-binding orbitals (NTBOs)
$\chi_a$~\cite{andersen84_prl53_2571,bloechl11_prb84_205101}. Inside the
augmentation sphere of radius \key{RAUG/RCOV} the head function of an
orbital is an augmented screened Hankel function, a combination of an AE
partial wave $\phi$ and its energy derivative $\dot\phi$,
\begin{equation}
  K^{\mathrm{aug}} = \phi\,k_\phi+\dot\phi\,k_{\dot\phi},
\end{equation}
matched to a Hankel function with $\kappa^2$ set by \key{K2}. In the
spheres of the neighbours it is continued by screened Bessel functions
$\bar J^{\mathrm{aug}}=\dot\phi\,j_{\dot\phi}$ with the structure constants
$\bar S$. Exponential tails (\key{TAILLAMBDA}, \key{RTAIL/RCOV}) and a
cutoff (\key{RTAILCUT/RCOV}) make the orbitals short ranged. The
orbitals considered per $\ell$ are selected by \key{NOFL}.

\section{Projection onto the local orbitals}\label{sec:projection}

The onsite density matrix
$D^\sigma_{ab}=\sum_n f_{n\sigma}\braket{\chi_a}{\psi_{n\sigma}}\braket{\psi_{n\sigma}}{\chi_b}$
requires the coefficients of the wave functions in the local-orbital
basis. Inside the spheres the one-center expansion
$\psi^1=\sum_i\phi_i\braket{\pt_i}{\tilde\psi}$ is fitted by local
orbitals in the least-squares sense,
\begin{equation}
  \bm c = \bigl[\bra{\chi}\Theta\ket{\chi}\bigr]^{-1}
          \bra{\chi}\Theta\ket{\phi}\braket{\pt}{\tilde\psi},
\end{equation}
where $\Theta$ restricts integrals to the spheres. Both overlaps are
computed with the exact $\chi$ inside the spheres, i.e.\ with
$K^{\mathrm{aug}}$ and $\bar J^{\mathrm{aug}}$ as radial functions.

\paragraph{Pitfall.} The original implementation represented $\chi$ by
its partial-wave expansion $\braket{\pt}{\tilde\chi}$. The $\dot\phi$
component of $\chi$ converges slowly in that expansion, so the fitted
orbital lacked norm and the density matrix overcounted charge: 13--15
electrons on Ni in NiO instead of about 7. The hybrid energy and the gap
then changed with the number of projectors (NiO \pbezero\ gap 3.96,
4.08, 4.10~eV for 2, 3, 4 $d$ projectors). With the exact overlaps they
are converged (2.99, 3.00, 3.00~eV). The old fit can be restored with the
environment variable \texttt{VESPAW\_NTBO\_OLDPROJ=1} for comparison.

\section{\pbezero}

\pbezero\ keeps only the onsite terms of the PBE0 exchange correction,
\begin{equation}
  E_{\mathrm{xc}}^{\pbezero} = E_{\mathrm{xc}}^{\mathrm{PBE}}[n]
  + a_x\bigl(E_x^{\mathrm{HF},\onec} - E_x^{\mathrm{PBE},\onec}\bigr),
  \qquad
  E_x^{\mathrm{HF},\onec} = -\frac12\sum_\sigma\sum_{abcd}
  D_{ac}^\sigma D_{bd}^\sigma U_{abcd},
  \label{eq:pbe0r}
\end{equation}
with the onsite Coulomb tensor $U$ of the local orbitals. The
double-counting term $E_x^{\mathrm{PBE},\onec}$ is the PBE exchange of the
local-orbital density, obtained from the full one-center density with the
weight $w(r)=(n_{\mathrm{loc}}/n_{\mathrm{all}})^2$. Core--valence exchange
can be included (\key{CV=T}). The mixing $a_x$ is set globally by
\key{!CONTROL!DFT!NTBO HFWEIGHT} and per species by \key{LHFWEIGHT}. The
exchange correction enters the Hamiltonian in two places: the one-center
PAW Hamiltonian $\Delta H$ (derivative with respect to the full one-center
density) and an operator in the local-orbital basis (derivative with
respect to $D$).

\section{\pbeslzero}\label{sec:pbesl0r}

\pbeslzero\ (called \scanzero\ before October 2026) replaces the semilocal
onsite exchange by SCAN-L exchange while keeping PBE for the delocalized
part and for all correlation:
\begin{align}
  E_{\mathrm{xc}}^{\pbeslzero} &= E_{\mathrm{xc}}^{\mathrm{PBE}}[n]
  - E_x^{\mathrm{PBE},\onec}
  + a_x E_x^{\mathrm{HF},\onec} + (1-a_x) E_x^{\mathrm{SCAN},\onec}
  \nonumber\\
  &= E_{\mathrm{xc}}^{\mathrm{PBE}}[n]
  + \underbrace{\bigl(E_x^{\mathrm{SCAN},\onec}-E_x^{\mathrm{PBE},\onec}\bigr)}_{\text{swap term}}
  + a_x\bigl(E_x^{\mathrm{HF},\onec} - E_x^{\mathrm{SCAN},\onec}\bigr).
  \label{eq:pbesl0r}
\end{align}
The swap term is not scaled by $a_x$. \pbeslzero\ is selected by
\key{!NTBO MODUS='HYBRID' HFWEIGHT=$a_x$ DCTYPE='SCANL'} with a PBE base
functional. A consistent implementation requires:
\begin{enumerate}
\item the swap term in both Hamiltonian channels (one-center and
  local-orbital), at full weight;
\item the full Euler--Lagrange potential of the onsite SCAN-L exchange,
  Eq.~\eqref{eq:euler}, with the weight $w$ inside the derivatives;
\item exact spin scaling, Eq.~\eqref{eq:spinscaling}. Evaluating the
  onsite SCAN-L exchange for the spin-summed density replaces
  spin-polarized PBE exchange by unpolarized SCAN-L exchange in the swap
  term and removes a fraction $1-a_x$ of the onsite exchange splitting;
  magnetic gaps collapse (NiO 1.0 instead of 2.0~eV);
\item the exact local-orbital projection of Sec.~\ref{sec:projection}.
\end{enumerate}
All four were found by checking that the variational optimizer and
damped dynamics reach the same state (Sec.~\ref{sec:consistency}).

\section{\scanzero}\label{sec:scan0r}

\scanzero\ uses SCAN with the orbital kinetic energy density everywhere,
for the plane-wave density and in the spheres, for exchange and
correlation, and replaces a fraction $a_x$ of the onsite SCAN exchange of
the local orbitals by exact exchange:
\begin{equation}
  E_{\mathrm{xc}}^{\scanzero} = E_{\mathrm{xc}}^{\mathrm{SCAN}}[n,\tau]
  + a_x\bigl(E_x^{\mathrm{HF},\onec} - E_x^{\mathrm{SCAN},\onec}\bigr).
  \label{eq:scan0r}
\end{equation}
It is the meta-GGA analogue of \pbezero, Eq.~\eqref{eq:pbe0r}, without a
swap term, and is selected by \key{TYPE='SCAN' TSCANDELOC=T} with
\key{!NTBO MODUS='HYBRID' HFWEIGHT=$a_x$}. The double counting
$E_x^{\mathrm{SCAN},\onec}$ uses the onsite $\tau$ of the local orbitals,
is spin scaled, and by default contains exchange only (\key{TSCANCORR=T}
adds SCAN correlation). Exchange-only spin scaling and the
spin-polarized correlation of Chapter~\ref{chap:dft} apply as for SCAN.
With \key{TYPE='R2SCAN'} the same construction gives r$^2$SCAN0r, with
r$^2$SCAN everywhere and in the double counting.


\section*{Further reading}
The original work on the subject of this chapter is found in
Refs.~\cite{adamo99_jcp110_6158,bloechl11_prb84_205101,sotoudeh17_prb95_235150}.

% =====================================================================
\chapter{Self-consistency: the variational optimizer}\label{chap:variational}
\minitoc
\section{Introduction}

The Kohn--Sham equations are a fixed-point problem: the density
determines the potential, and the potential determines the density. In
almost every plane-wave code this problem is solved by mixing input and
output densities until they agree, a procedure that works well in most
cases and fails, when it fails, without warning. The theory, however,
offers a minimum principle. This chapter shows how the self-consistency
problem can be written as a minimization in which the effective
Hamiltonian is the variable, and every energy is that of an actual
state. It describes the structure of this Hamiltonian space, including
a gauge mode that is specific to the PAW method, and the line search
that turns it into a robust algorithm.

% =====================================================================

\section{The problem with density mixing}

The standard SCF cycle alternates between solving the Kohn--Sham
equations in a fixed potential and building a new input potential by
mixing the input and output densities or potentials (Pulay/DIIS
\cite{Pulay1980}, Anderson~\cite{Anderson1965}, with Kerker
preconditioning~\cite{Kerker1981}). The energy along this path is a
Harris--Foulkes-type estimate~\cite{Harris1985,FoulkesHaydock1989}; it is
not an upper bound and need not decrease. Mixing can stagnate, oscillate,
or converge to an excited state, and nothing in the iteration signals
it. VESPAW offers density mixing (\key{MIXER='PULAY'} or
\key{'ANDERSON'}) but uses a variational optimizer by default.

\section{Free energy as a function of the Hamiltonian}\label{sec:route1}

We minimize the Mermin free energy, Eq.~\eqref{eq:mermin}, in the PAW
form. The parameters of the Hamiltonian are collected in a vector
\begin{equation}
  \vth = \bigl(\tilde v(\bm r),\; v_\tau(\bm r),\; \{\Delta H_{R,ij}\}\bigr).
\end{equation}
For given $\vth$ we determine the eigenstates $\psi_n(\vth)$ of
$\tilde H[\vth]$ and the Fermi occupations $f_n(\vth)$ at fixed $N$, and
define
\begin{equation}
  \mathcal{A}(\vth) = A\bigl[\{\psi_n(\vth)\},\{f_n(\vth)\}\bigr].
\end{equation}
Because $\psi_n(\vth)$ and $f_n(\vth)$ form an admissible ensemble,
$\mathcal{A}(\vth)$ is an upper bound of the minimum of $A$ for every
$\vth$. At self-consistency, where $\vth$ equals the parameters
$\vout[\vth]$ built from the output state, the bound is attained.
Following Gonze~\cite{Gonze1996},
\begin{equation}
  \frac{\partial \mathcal{A}}{\partial \vth}
  = \chi(\vth)\,\bigl(\vout[\vth]-\vth\bigr),
\end{equation}
with $\chi$ the static independent-particle response. Since $\chi$ is
negative semidefinite, the residual $\bm r=\vout-\vth$, and any
positive-definite preconditioned version $K\bm r$, is a descent
direction; $\chi$ never has to be built. The Kerker filter
$K(G)=G^2/(G^2+q_0^2)$ is applied to the plane-wave part of the residual,
with $q_0^2$ set by \key{MIXG0SQ} (default 1.5~bohr$^{-2}$). The default is
made for solids. For molecules and slabs with large vacuum it damps the
long-wavelength components so strongly that the gains of a step fall
below the noise of the iterative eigensolver, and the line search
stalls; use $q_0^2\approx0.05$--0.1 there.

\section{The gauge direction}

A uniform shift $\tilde H\to\tilde H+c\hat O$ shifts all eigenvalues and
the Fermi level by $c$ and leaves orbitals and occupations unchanged. In
parameter space it is the direction $\bm s=(1,0,\{\Delta O_{R,ij}\})$, and
$\mathcal A(\vth+c\bm s)=\mathcal A(\vth)$ exactly. The two halves of
$\bm s$ are updated by different rules (Kerker on the grid, damped steps
in $\Delta H$). If $\bm s$ is not treated explicitly, they drift apart and
act as a spurious potential step between spheres and interstitial (Si:
convergence to an excited state 28~mH above the ground state). VESPAW
therefore
\begin{enumerate}
\item projects $\bm s$ out of every residual, and
\item after each step moves $\vth$ along $\bm s$ so that its mean grid
  potential equals that of the output state.
\end{enumerate}
The second step does not change any energy, but without it $\vth$ keeps
its starting gauge while the true potential moves: eigenvalues and Fermi
level are then offset (0.49~eV in Si), and with them the Lagrange
multipliers in forces and stress (Chapter~\ref{chap:forces}).

\section{Line search and acceptance}

Steps are proposed by an Anderson extrapolation built from
\emph{accepted} points only. A trial $\vth'$ is accepted if
$\mathcal A(\vth')\le\mathcal A_{\mathrm{best}}+\epsilon$; otherwise the
step is halved towards the best point. After five consecutive rejections
the best point is re-evaluated (``re-anchoring''), because its energy may
stem from an eigensolution that was not fully converged.

\begin{algorithm}[t]
\caption{Hamiltonian-space variational SCF (Route 1).}
\begin{algorithmic}[1]
\State $\vth_0\gets\vout[\psi_{\mathrm{start}}]$; hold fixed for a short eigensolver warm-up
\Loop
  \State diagonalize $\tilde H[\vth]$ (Davidson, RMM-DIIS or CG); Fermi occupations at fixed $N$
  \State evaluate $\mathcal A(\vth)$ at the output state; build $\vout$
  \State $\bm r\gets\vout-\vth$; project out $\bm s$
  \If{$\mathcal A(\vth)\le\mathcal A_{\mathrm{best}}+\epsilon$}
    \State accept; add $(\vth,K\bm r)$ to the history; stop if converged
    \State $\vth\gets$ Anderson proposal
  \Else
    \State $\vth\gets\frac12(\vth+\vth_{\mathrm{best}})$ \Comment{re-anchor after 5 rejections}
  \EndIf
  \State re-gauge $\vth$ along $\bm s$ to the mean grid potential of $\vout$
\EndLoop
\end{algorithmic}
\end{algorithm}

Convergence is reached when the mean-square grid residual is below
\key{MIXTOL} (and the one-center residual below its square root), or when
two consecutive accepted energies differ by less than \key{EDIFF}. Route 1
is selected by \key{!PSIDYN MIXER='VARIATIONAL'} together with
\key{OPTIMIZER='DAV'}, \key{'RMM'} or \key{'CG'}; it is the default mixer.

\paragraph{Hybrid operators.} The onsite hybrid operator of
Chapter~\ref{chap:hybrids} depends on the orbitals through $D$, not only
on the density. It is taken from the current state rather than
line-searched; including its matrix elements in $\vth$ left a residual
that did not decay.

\section{Route 2: direct joint minimization}

As an alternative (\key{OPTIMIZER='EDFT'}), $A$ is minimized directly with
respect to the orbitals and an auxiliary Hamiltonian $H^{\mathrm{aux}}$
whose eigenvalues define the occupations, in the spirit of ensemble
DFT~\cite{MarzariVanderbiltPayne1997,FreysoldtBoeckNeugebauer2009}. The
orbital direction is the Teter-preconditioned~\cite{TeterPayneAllan1989}
projected residual combined by Polak--Ribi\`ere conjugate gradients; the
direction for $H^{\mathrm{aux}}$ is $H^{\mathrm{phys}}-H^{\mathrm{aux}}$.
Trial points are built by L\"owdin orthonormalization and rotated into
the eigenbasis of $H^{\mathrm{aux}}$; step lengths come from a quadratic
model, with the same acceptance and re-anchoring logic as Route~1. The
orbital direction is computed from a Kerker-preconditioned Hamiltonian.
With this preconditioning Route~2 currently stops about $5\times10^{-6}$~H
above the minimum on Si; without it (\texttt{VESPAW\_EDFTVAR\_NOKERKER=1}) it
reaches the minimum.

\section{The consistency test}\label{sec:consistency}

The variational optimizer minimizes the energy. Damped Car--Parrinello
dynamics (Chapter~\ref{chap:eigensolvers}) and density mixing converge to
a self-consistent state. The two coincide only if the Hamiltonian is the
derivative of the energy. Agreement between them, to $10^{-6}$~H or
better, is therefore a sharp test of any new functional or
implementation: a missing potential term shows up as the variational
optimizer finding energies \emph{below} the self-consistent state without
converging. The regression suite (Chapter~\ref{chap:code}) applies this
test to every case.


\section*{Further reading}
The original work on the subject of this chapter is found in
Refs.~\cite{Gonze1996,MarzariVanderbiltPayne1997,FreysoldtBoeckNeugebauer2009,VariationalDraft}.

% =====================================================================
\chapter{Eigensolvers and dynamics}\label{chap:eigensolvers}
\minitoc
\section{Introduction}

Most of the computing time of a calculation is spent in finding the
lowest eigenstates of the Hamiltonian. For a few hundred bands in a
basis of tens of thousands of plane waves, a direct diagonalization is
out of the question. Iterative methods instead improve a set of trial
wave functions, using only the action of the Hamiltonian on them. In
the original Car--Parrinello approach this is done by a fictitious
dynamics of the wave functions. This chapter describes the iterative
eigensolvers of VESPAW and the damped dynamics of CP-PAW, and what each
of them requires to return reliable, orthonormal eigenstates.

% =====================================================================

For Route~1 the eigenstates of $\tilde H[\vth]$ must be found at every
iteration. They are warm-started from the previous iteration and
converged iteratively. All solvers work with the generalized problem
$\tilde H\psi=\epsilon\hat O\psi$ and apply $\tilde H$ and $\hat O$ in
blocks of bands.

\section{Block Davidson}

\key{OPTIMIZER='DAV'}. Each band is extended by a preconditioned
correction vector $(\tilde H-\epsilon\hat O)\psi$; the new vectors are
projected out of \emph{all} current bands (twice), vectors lost to
cancellation are dropped, and the extended basis is canonically
orthogonalized with an eigenvalue cutoff of $10^{-10}$. A Rayleigh--Ritz
step in this subspace gives the new bands. Converged bands are locked.
The implementation keeps seven wave-function-sized arrays. Davidson is
robust from any start and is the reference solver.

\paragraph{Pitfall.} An earlier Davidson orthogonalized only against
converged bands and stabilized the near-singular overlap with a shift;
the returned orbitals were not orthonormal, which produced energies below
the minimum and split degenerate states.

\section{RMM-DIIS}

\key{OPTIMIZER='RMM'}. Residual minimization with direct inversion in the
iterative subspace~\cite{WoodZunger1985,KresseFurthmuller1996PRB}
minimizes the norm of the residual of each band separately, without
orthogonalization against other bands, and is therefore cheap and
parallel. It converges to the eigenvector nearest the starting guess.
Started from random orbitals it converges to wrong states, so the first
\key{RMMDELAY} iterations (default 8) use Davidson. Bands are processed
in blocks of 32 (\key{NRMMX} inner sweeps). RMM-DIIS is the fastest
option for large systems (Si$_{64}$: 247~s vs 386~s for Davidson).

\paragraph{Pitfall.} With the onsite hybrid, RMM-DIIS reaches the correct
energy but can miss the lowest empty band (MgO \pbeslzero\ gap 6.54 instead
of 4.23~eV). Take band gaps of hybrid calculations from Davidson.

\section{Conjugate gradients}

\key{OPTIMIZER='CG'}. Band-by-band preconditioned conjugate-gradient
minimization of the Rayleigh quotient with orthogonalization to the
lower bands, followed by a subspace rotation. Reliable, slower than the
block methods for many bands.

\section{Car--Parrinello dynamics with friction}

The original CP-PAW optimizer (\key{OPTIMIZER='MD'}, the default when no
eigensolver is chosen) treats the wave-function coefficients as classical
degrees of freedom with a fictitious mass~\cite{CarParrinello1985},
\begin{equation}
  m_\psi\ddot{\tilde\psi}_n=-\tilde H\tilde\psi_n
  +\sum_m\hat O\tilde\psi_m\Lambda_{mn}-\gamma\,m_\psi\dot{\tilde\psi}_n,
\end{equation}
with Lagrange multipliers $\Lambda$ that keep the orbitals orthonormal and
a friction $\gamma$ (\key{!PSIDYN FRIC}). With friction the dynamics
relaxes into the ground state. The ``autopilot'' (\key{!PSIDYN!AUTO})
lowers the friction when the energy decreases and raises it when it
increases. Damped dynamics is slow but very robust, converges to the
self-consistent state for any functional, and serves as the reference
in the consistency test. \key{SAFEORTHO=T} orthogonalizes such that the
energy is conserved; it should not be used with variable occupations.


\section*{Further reading}
The original work on the subject of this chapter is found in
Refs.~\cite{Davidson1975,WoodZunger1985,KresseFurthmuller1996PRB,Payne1992,CarParrinello1985}.

% =====================================================================
\chapter{Forces, stress and molecular dynamics}\label{chap:forces}
\minitoc
\section{Introduction}

Once the electrons are in their ground state, the atoms feel forces and
the cell feels a stress. These derivatives of the energy are what makes
structure optimization and molecular dynamics possible. In a
variational theory they follow from the explicit dependence of the
energy on the atomic positions, but in the PAW method this dependence
hides in many places: the projector functions, the compensation charges
and the overlap operator all move with the atoms. This chapter collects
the expressions for forces and stress and describes the equations of
motion for the atoms.

% =====================================================================

\section{Forces}

The force on atom $R$ is the total derivative of the energy with respect
to its position. At a stationary point of the energy the orbital
response drops out (Hellmann--Feynman~\cite{Feynman1939}). In PAW the
projectors, the compensation charges and the one-center terms move with
the atom, and the orthonormality constraint depends on the positions
through $\hat O$:
\begin{equation}
  \bm F_R=-\frac{\partial E}{\partial\bm R}\Big|_{\mathrm{explicit}}
  -\sum_{nm}\Lambda_{mn}\,\bra{\tilde\psi_n}\frac{\partial\hat O}{\partial\bm R}\ket{\tilde\psi_m}
  \times(\text{occupation weights}),
\end{equation}
where the first term contains the derivatives of $\Delta H$, the projectors
and the compensation charges. In Car--Parrinello dynamics $\Lambda$ comes
from the orthonormality constraint. For the eigensolvers the orbitals are
Ritz vectors, and $\Lambda=\mathrm{diag}(\epsilon_n)$.

\paragraph{Pitfalls.} Before September 2026 the eigensolvers did not set
$\Lambda$ at all, so the constraint term was missing (0.5\% error in the
Si force), and with the variational mixer the eigenvalues carried the
gauge offset of the previous chapter. Both are now fixed, and forces
agree with finite differences of the energy to $10^{-4}$ relative
(regression case \texttt{si\_forces}).

\section{Stress}

The stress tensor is the derivative of the energy with respect to a
homogeneous strain $\epsilon$,
\begin{equation}
  \sigma_{\alpha\beta}=-\frac1V\frac{\partial E}{\partial\epsilon_{\alpha\beta}},
\end{equation}
following Nielsen and Martin~\cite{NielsenMartin1985}; positive values
mean that the cell wants to expand, and the pressure is
$\mathrm{tr}\,\sigma/3$. It contains kinetic, Hartree, exchange-correlation,
projector and compensation-charge terms, and the same constraint term as
the forces.

The analytic stress is taken at a fixed set of plane waves (the basis
strains with the cell). A finite difference at fixed cutoff changes the
basis and contains the \emph{Pulay stress}, which vanishes as the cutoff
grows (Si: 2\% at 25~Ry, 0.5\% at 50~Ry). Lattice constants from
energy-volume curves need converged cutoffs for the same reason.

\paragraph{Output.} \key{!CONTROL!GENERIC FORCES=T STRESS=T} evaluates
forces and stress also for fixed atoms and a fixed cell and prints them at
full precision at the end of the run (\texttt{FORCE}, \texttt{STRESS},
\texttt{PRESSURE[GPA]} lines).

\section{Structure relaxation and molecular dynamics}

Atomic motion is controlled by \key{!CONTROL!RDYN}. With friction
(\key{FRIC}, \key{!RDYN!AUTO}) the atoms relax into a minimum; without
friction they follow Newton's equations. The cell shape can be relaxed or
propagated with \key{!CONTROL!CELL MOVE=T}~\cite{NielsenMartin1985}.

\paragraph{Car--Parrinello vs Born--Oppenheimer MD.} In Car--Parrinello
MD the wave functions are propagated together with the atoms and stay
close to the ground state if their fictitious mass is small enough
(adiabaticity~\cite{BloechlParrinello1992}). The time step must resolve
the fastest wave-function motion. In Born--Oppenheimer MD the electrons
are converged at each step (with an eigensolver and Route~1); the time
step is set by the ions, but each step costs a full SCF.

\paragraph{Thermostats.} Temperature is controlled by Nos\'e--Hoover
thermostats~\cite{Nose1984} on the atoms (\key{!RDYN!THERMOSTAT T[K]}) and,
in Car--Parrinello MD, on the wave functions~\cite{Bloechl2002}
(\key{!PSIDYN!THERMOSTAT}), which removes the heat that leaks from the
ions into the fictitious electronic motion.


\section*{Further reading}
The original work on the subject of this chapter is found in
Refs.~\cite{NielsenMartin1985,Nose1984,BloechlParrinello1992,Bloechl2002}.

% =====================================================================
\chapter{Structure of the code}\label{chap:code}
\minitoc
\section{Introduction}

This chapter is written for those who want to read the source code,
change it, or add to it. It gives an overview of the source files,
follows one self-consistency step through the program, and describes
how a new functional or a new optimizer is added and tested.

% =====================================================================

\section{Source files}

\begin{center}
\begin{tabular}{@{}ll@{}}
\toprule
File & Content \\
\midrule
\texttt{vespaw.f90}, \texttt{vespaw\_driver.f90} & main program, time-step loop, reports \\
\texttt{vespaw\_ioroutines.f90} & reading \texttt{.cntl}/\texttt{.strc} input \\
\texttt{vespaw\_waves1.f90}, \texttt{vespaw\_waves2.f90} & wave functions, density, $H\psi$, forces, total energy \\
\texttt{vespaw\_cg.f90} & Davidson, RMM-DIIS, CG \\
\texttt{vespaw\_mixer.f90} & variational, Pulay and Anderson mixers \\
\texttt{vespaw\_edft.f90} & Route~2 (EDFT) \\
\texttt{vespaw\_potential.f90} & plane-wave potential, electrostatics, XC on the grid \\
\texttt{vespaw\_augmentation.f90} & one-center (PAW sphere) terms \\
\texttt{vespaw\_setups.f90} & setup construction and tests \\
\texttt{vespaw\_lmto.f90} & local orbitals, \pbezero, \pbeslzero, \scanzero \\
\texttt{vespaw\_scantau.f90}, \texttt{vespaw\_scannsph.f90} & SCAN with orbital $\tau$, non-spherical one-center SCAN \\
\texttt{vespaw\_dft.f90} & exchange-correlation functionals \\
\texttt{vespaw\_atoms.f90}, \texttt{vespaw\_cell.f90} & ionic and cell dynamics \\
\texttt{vespaw\_occupations.f90} & occupations, Mermin \\
\bottomrule
\end{tabular}
\end{center}
The code is organized in CP-PAW's object style: each module keeps its
data in a Fortran module and is accessed through \texttt{OBJECT\$SET},
\texttt{\$GET} and action routines.

\section{One SCF step}

\begin{enumerate}
\item \texttt{PAW} (driver) calls \texttt{TIMESTEP}, which calls
  \texttt{WAVES\$ETOT}.
\item \texttt{WAVES\$ETOT}: projections $\braket{\pt}{\tilde\psi}$, density
  and one-center density matrices; plane-wave potential
  (\texttt{WAVES\_VOFRHO}); one-center potentials and $\Delta H$
  (\texttt{AUGMENTATION}); hybrid terms (\texttt{LMTO}); forces and
  stress (\texttt{WAVES\$FORCE}).
\item Variational mixer (\texttt{WAVES\_VARMIXHAMILTONIAN}): the output
  Hamiltonian and the energy are handed to the line search, which returns
  $\vth$ for the next diagonalization.
\item Eigensolver (\texttt{DAVIDSON\$STATE\_BY\_BLOCK},
  \texttt{RMM\$STATE\_BY\_BLOCK}, \texttt{CG\$STATE\_BY\_STATE}) or, for
  dynamics, $H\psi$ and propagation.
\item Reports (\texttt{PRINFO}), capture of forces and stress, and the
  switch of all dynamical variables to the next step.
\end{enumerate}

\section{Input}

The \texttt{.cntl} and \texttt{.strc} files are read into linked lists
(\texttt{LINKEDLIST\$...}). Blocks start with \texttt{!NAME} and end with
\texttt{!END}; a block whose name is changed (e.g.\
\texttt{!ORBPOT\_OFF}) is ignored. \texttt{READ=F} inside a block also
switches the block off. The routines \texttt{READIN\_*} in
\texttt{vespaw\_ioroutines.f90} transfer the values to the modules.

\section{Adding a functional}

\begin{enumerate}
\item Implement energy density and potential in \texttt{vespaw\_dft.f90}
  (or a module of its own), register a type name in
  \texttt{READIN\_DFTTYPEFROMNAME}.
\item Provide the one-center evaluation (\texttt{vespaw\_augmentation.f90}),
  with the full potential, Eq.~\eqref{eq:euler}.
\item Verify: energy vs finite difference of the potential along a
  density change; variational optimizer vs damped dynamics
  (Sec.~\ref{sec:consistency}); forces and stress vs finite differences.
\item Add a regression case.
\end{enumerate}

\section{Adding an optimizer}

Optimizers are dispatched on \texttt{OPTIMIZERTYPE} in \texttt{WAVES\$ETOT}.
A new eigensolver must return orthonormal Ritz vectors and their Ritz
values in \texttt{THIS\%EIGVAL} (used for the Lagrange multipliers of the
forces). It should be tested against Davidson on the whole regression
suite.

\section{Regression suite}

\texttt{regression/run.py} runs all cases (\texttt{regression/*/case.json})
or selected ones (\texttt{run.py CASE:RUN}). A case has template inputs
with \texttt{@PLACEHOLDER@} substitutions, several runs (optimizers,
cold and seeded starts, multi-stage runs sharing a restart file), checks
against reference values, agreement tests between runs, finite-difference
tests of forces and stress, and a list of known issues. The cases cover
PBE Si (all optimizers), hot Al (metal, Mermin), Si$_8$ at $\Gamma$, NiO
AFM (PBE, PBE+$U$, \pbezero, \pbeslzero, \scanzero, TSCAN), MgO
\pbeslzero, TSCAN Si, spin--orbit Si, a polar slab with dipole
correction, space-group symmetry, and forces and stress. Every change to the code should leave the suite passing.

\backmatter
\bibliographystyle{unsrtnat}
\bibliography{theory}
\end{document}
